%-> 封面与声明
\maketitle
\MAKETITLE
\makedeclaration

%-
%-> 中文摘要
%-
\intobmk\chapter*{\texorpdfstring{摘\quad 要}{摘要}}% 显示在书签但不显示在目录
\setcounter{page}{1}% 开始页码
\pagenumbering{Roman}% 页码符号

面向国产处理器平台完善 Android 运行时，是支撑 Android 应用生态跨架构迁移的重要基础。Android Runtime（ART，Android 运行时）承担 DEX 字节码解释、即时编译、提前编译、Java 本地接口桥接和运行时快速路径等核心职责，其执行引擎后端的完备度直接影响应用可运行性、启动开销与运行时吞吐。针对 Android Open Source Project（AOSP）15 官方版本尚未提供 LoongArch 架构 ART 执行引擎后端的问题，本文围绕“如何建立可运行且可验证的 LoongArch 平台 ART 执行环境，以及如何在统一实验条件下识别并优化平台相关性能瓶颈”这一核心问题展开研究。

本文首先完善了 LoongArch 平台 ART 执行引擎的关键模块，涵盖汇编器、C++ 解释器、Nterp、即时编译（Just-In-Time，JIT）编译器后端、提前编译（Ahead-Of-Time，AOT）与编译执行路径、Java 本地接口（Java Native Interface，JNI）桥接以及运行时快速路径，从而构建了集解释执行、编译执行、跨语言调用与运行时支撑于一体的完整执行链路。在此基础上，本文搭建了统一的基准测试评估框架，为后续的性能优化提供了标准、可复现的实验环境。

围绕统一实验框架，本文进一步从局部热点中提炼出 JIT 进入时机、运行时热路径、启动期校验链和后端向量化四类代表性问题，并设计和评估了自适应基线 JIT 提前触发机制、运行时热路径内存批量搬移机制、DEX 校验阶段化精简机制和 LSX 向量化目标指令映射（Lowering）路径。实验结果表明，自适应基线 JIT 机制在启动期短生命周期负载上取得可量化性能提升，其中 \texttt{startup.crypto.aes}、\texttt{startup.crypto.rsa} 与 \texttt{startup.crypto.signverify} 分别取得 \texttt{34.29\%}、\texttt{21.48\%} 与 \texttt{19.02\%} 的相对基线性能提升，并在 \texttt{SPECjvm2008/mpegaudio} 上将平均吞吐量由 \texttt{22.13\,ops/m} 提高到 \texttt{25.44\,ops/m}；运行时批量搬移机制在 \texttt{DaCapo/pmd} 上将运行时延由 \texttt{3272\,msec} 降至 \texttt{3244\,msec}，在 \texttt{mpegaudio} 上将吞吐量由 \texttt{25.38\,ops/m} 提高到 \texttt{25.56\,ops/m}。DEX 校验阶段化精简在 \texttt{startup.helloworld} 上取得最高 \texttt{5.46\%} 的启动期性能提升；补齐 LoongArch64 最小可用 LSX 浮点向量 Lowering 链路后，\texttt{SPECjvm2008/scimark.lu.small} 相较于对照版本取得 \texttt{44.40\%} 的吞吐量提升，并在重复验证中得到约 \texttt{45.70\%} 的平均性能提升。

研究结果表明，AOSP 15 ART 在 LoongArch 平台上的执行引擎适配与性能优化具有工程可行性；面向新架构平台的 ART 性能结论必须建立在统一验证框架、固定对照版本和明确工作负载边界之上；同时，LoongArch 平台 ART 的优化效果具有路径相关性和工作负载相关性，局部机制的有效性不能简单外推为全局叠加性能提升。

\keywords{LoongArch，Android Runtime，执行引擎适配，JIT 编译，性能优化}% 中文关键词
%-
%-> 英文摘要
%-
\intobmk\chapter*{Abstract}% 显示在书签但不显示在目录

Improving Android runtime support on domestically developed processor platforms is an important foundation for supporting cross-architecture migration of the Android application ecosystem. Android Runtime (ART) is responsible for core tasks including DEX bytecode interpretation, Just-In-Time compilation, Ahead-Of-Time compilation, Java Native Interface bridging, and runtime fast paths. The completeness of its execution engine backend directly affects application runnability, startup overhead, and runtime throughput. Since Android Open Source Project (AOSP) 15 does not yet provide an official ART execution engine backend for LoongArch, this thesis focuses on the central problem of how to establish a runnable and verifiable ART execution environment on LoongArch and how to identify and optimize platform-related performance bottlenecks under unified experimental conditions.

This thesis first completes the key modules of the ART execution engine on LoongArch, including the assembler, the C++ interpreter, Nterp, the Just-In-Time (JIT) compiler backend, Ahead-Of-Time (AOT) compilation and compiled execution paths, Java Native Interface (JNI) bridging, and runtime fast paths. These modules build a complete execution chain integrating interpreted execution, compiled execution, cross-language calls, and runtime support. On this basis, this thesis builds a unified benchmark evaluation framework, providing a standard and reproducible experimental environment for subsequent performance optimization.

Around the unified experimental framework, this thesis further extracts four representative problems from local hot spots: JIT entry timing, runtime hot paths, startup-time verification chains, and backend vectorization. It then designs and evaluates an adaptive baseline JIT early triggering mechanism, a batched memory movement mechanism for runtime hot paths, a staged simplification mechanism for DEX verification, and an LSX vectorization target-instruction lowering path. The experimental results show that the adaptive baseline JIT mechanism achieves quantifiable performance improvements on short-lived startup workloads, with relative improvements over the baseline of \texttt{34.29\%}, \texttt{21.48\%}, and \texttt{19.02\%} on \texttt{startup.crypto.aes}, \texttt{startup.crypto.rsa}, and \texttt{startup.crypto.signverify}, respectively, and increases the average throughput of \texttt{SPECjvm2008/mpegaudio} from \texttt{22.13\,ops/m} to \texttt{25.44\,ops/m}. The runtime batched memory movement mechanism reduces the runtime latency of \texttt{DaCapo/pmd} from \texttt{3272\,msec} to \texttt{3244\,msec} and increases the throughput of \texttt{mpegaudio} from \texttt{25.38\,ops/m} to \texttt{25.56\,ops/m}. The staged simplification of DEX verification achieves a maximum startup-period performance improvement of \texttt{5.46\%} on \texttt{startup.helloworld}. After completing the minimal usable LoongArch64 LSX floating-point vector lowering chain, \texttt{SPECjvm2008/scimark.lu.small} obtains a \texttt{44.40\%} throughput improvement over the control version, and repeated validation shows an average performance improvement of approximately \texttt{45.70\%}.

The research results show that adapting and optimizing the AOSP 15 ART execution engine on LoongArch is technically feasible. They also show that ART performance conclusions for a new architecture must be established on a unified validation framework, a fixed control version, and explicit workload boundaries. Meanwhile, optimization effects on LoongArch ART are path-dependent and workload-dependent, so the effectiveness of local mechanisms cannot be simply extrapolated into a global cumulative performance improvement.

\KEYWORDS{LoongArch, Android Runtime, execution engine adaptation, JIT compilation, performance optimization}% 英文关键词
